\section{引言}\label{sect_introduction1}%
%

QCD 的格点计算已经展示了与大量实验测量互补的能力，在某些情况下其精度甚至超过实验测量。如今，部分子分布函数（PDF）的格点计算已提上日程。除了允许计算 PDF 前几个 Mellin 矩的标准途径之外，人们正在探索直接在动量分数空间获取 PDF 的新技术。现有的各种方案~\cite{Aglietti:1998ur,Abada:2001if,Detmold:2005gg,Braun:2007wv,Ji:2013dva,Ma:2017pxb}细节各异，但具有一个共同的一般框架：利用连续理论中的 QCD 共线因子化，从适当的欧氏关联函数的格点计算中提取 PDF。

一个特别流行的建议~\cite{Ji:2013dva}引发了近来的大量活动 \cite{Lin:2014zya,Alexandrou:2015rja,Chen:2016utp,Alexandrou:2016jqi,Lin:2017ani,Zhang:2017bzy,Chen:2017gck,Chen:2017mzz,Alexandrou:2017huk, Alexandrou:2018pbm,Chen:2018xof,Chen:2018fwa,Alexandrou:2018eet,
Lin:2018qky}，它引入了部分子\emph{准分布}（quasi-distribution，qPDF）的概念，定义为以 Wilson 线相连的非定域夸克—反夸克算符的 Fourier 变换。qPDF 随后通过匹配得到 PDF：既可直接在 $\MSbar$ 方案中匹配，也可经由大动量因子化方案这一中间步骤（``大动量有效理论''（LaMET）~\cite{Ji:2014gla,Izubuchi:2018srq}）。后一技术有助于强调：经 Fourier 变换之后，强子动量 $p$ 仍是唯一的有量纲参量。因此，QCD 耦合的相关能标在差一个无量纲常数的意义上由 $p$ 决定，而且高扭度修正一般受动量 $p$ 的幂次压制。一个密切相关的量——\emph{赝分布}（pseudo-distribution，pPDF）则在文献~\cite{Radyushkin:2017cyf,Orginos:2017kos,Karpie:2017bzm}中引入。

使用这些新方法的 PDF 格点计算目前正从探索阶段迈向精确计算阶段，因此必须回答诸如高扭度（幂次压低）修正是否得到很好控制之类的问题。研究高扭度影响的一种可能途径是，把这类修正作为自由参数引入，从多个欧氏关联函数的整体分析中提取 PDF。该做法在文献 \cite{Bali:2018spj} 中以 $\pi$ 介子光锥分布振幅（LCDA）为例并采用位置空间策略~\cite{Braun:2007wv,Bali:2017gfr}得到了展示。
%
在这一情形中，高扭度效应的解析结构是清楚的，因此可以只用一个自由参数来建模。
一般情形则可能更复杂。考虑到格点计算中可采用的观测量众多（见例如~\cite{Ma:2017pxb}），
最好能有一个一般性方法来估计相应的高扭度修正及其预期的 Bjorken-$x$ 依赖。
本文的目的在于指出：此类观测量中幂次压低贡献的问题可以借助 renormalon 概念 \cite{Beneke:1998ui,Beneke:2000kc}来处理。下文我们将这一技术
应用于 qPDF 与 pPDF。

用 renormalon 研究幂修正的方法基于如下事实：由于耦合常数跑动破坏了 QCD 的标度不变性，不同扭度的算符在重正化下相互混合。在截断方案中这种混合是显式的；而在维数正规化中，它表现为微扰级数的阶乘发散。物理观测量对因子化尺标的无关性意味着不同扭度之间存在错综复杂的相消——即 renormalon 模糊性的相消。反过来，领头扭度表达式中这些模糊性的存在可用来估计幂次压低修正的大小。在概念上，这类似于通过对数尺度依赖来估计固定阶微扰结果的精度。renormalon 方法此前已被用于研究深度非弹性散射中高扭度修正的 Bjorken-$x$ 依赖 \cite{Beneke:1995pq,Dokshitzer:1995qm,Dasgupta:1996hh}、碎裂函数~\cite{Dasgupta:1996ki,Beneke:1997sr}、$\pi$ 介子 LCDA~\cite{Braun:2004bu}以及横动量依赖的部分子分布~\cite{Scimemi:2016ffw}。

为了说明 renormalon 概念如何用于洞察幂修正的结构，让我们考虑准分布的常用表达式~\cite{Ji:2015qla,Izubuchi:2018srq}
\begin{align}
\label{qPDF1}
 \mathcal{Q}(x,p) &= \int_{-1}^1 \!\frac{dy}{|y|} C_{\Qd}(\tfrac{x}{y},xp,\mu_F) q(y,\mu_F) + \frac{1}{p^2}  \mathcal{Q}_4(x,p) + \ldots
\end{align}
其中 $q(y,\mu_F)$ 是夸克 PDF，$p$ 是强子动量，$x$ 为动量分数。为简洁起见我们不再显式写出对重正化尺度的依赖。为避免大对数，因子化尺度 $\mu_F$ 必须取为 $|x| p$ 的量级。系数函数 $C(x,p,\mu_F) = \delta(1-x) + \mathcal{O}(\alpha_s)$ 由微扰展开给出。${\cal O}(\alpha_s)$ 修正在味道非单态情形首先由文献 \cite{Xiong:2013bka} 计算，另见~\cite{Stewart:2017tvs}。

要理解 renormalon 的作用，必须仔细审视式~\eqref{qPDF1} 中领头项与高扭度附加项之间的划分。让我们暂时假定因子化是用硬截断 $\Lambda_{\rm QCD} \ll \mu_F  \ll p$ 完成的，也就是说圈动量 $|k| > \mu_F$ 的贡献归入系数函数，而 $|k| < \mu_F$ 的贡献归入 PDF。在此方案中，当 $p\to\infty$ 时系数函数有如下展开：
\begin{align}
\label{C-cutoff}
 C_{\Qd}(x,p,\mu_F) &= \delta(1-x) + c_1 \alpha_s + c_2\alpha_s^2 + \ldots
\notag\\ &\quad{}
- \frac{\mu_F^2}{p^2} D_{\Qd}(x) +\ldots\,,
\end{align}
其中 $c_k= c_k(x,\ln p^2/\mu_F^2)$ 是对尺度仅有对数依赖的微扰系数，而 $D_{\Qd}$ 项代表领头幂修正。由于式~\eqref{qPDF1} 左端不依赖于 $\mu_F$，右端的任何此类依赖都应当相消。特别地，$c_k(x,\ln p^2/\mu_F^2)$ 中的对数尺度依赖被 PDF $q(x,\mu_F)$ 的尺度依赖所抵消。另一方面，幂次依赖的相消必然涉及扭度四贡献 $\mathcal{Q}_4(x,p)$。于是，在这种因子化方案中人们预期
\begin{align}
\label{Q4-cutoff}
 \mathcal{Q}_4(x,p,\mu_F) = \mu_F^2 \int_{-1}^1 \!\frac{dy}{|y|} D_{\Qd}(\tfrac{x}{y}) q(y,\mu_F) + \widetilde{\mathcal{Q}}_4(x,p,\mu_F)\,,
\end{align}
其中 $\widetilde{\mathcal{Q}}_4$ 对 $\mu_F$ 至多只有对数依赖。
项 $\sim \mu_F^2$ 的出现可追溯到截断方案中负责幂修正的扭度四算符的二次紫外发散（除对数紫外发散之外）。
可以证明，高扭度算符的截断依赖 $\sim \mu_F^2$ 正是式~\eqref{C-cutoff} 所给出的形式。

实践中，微扰计算通常采用维数正规化。此时式~\eqref{C-cutoff} 那样的幂次项不会出现。代价是在 $\MSbar$ 方案中计算的系数 $c_k$ 随阶数 $k$ 阶乘增长。阶乘增长意味着微扰级数之和仅在幂次精度内有定义，这一模糊性（renormalon 模糊性）必须由添加非微扰的高扭度修正来补偿。细致的分析表明~\cite{Beneke:1998ui,Beneke:2000kc}，系数的大 $k$ 发散行为（即 renormalon）与对极端（很小或很大）圈动量的敏感性一一对应。特别地，领头扭度系数函数中的红外 renormalon 被扭度四算符矩阵元中的紫外 renormalon 所补偿。这样一来，与截断方案相同的图像在维数正规化中重新出现。

回到 \eqref{Q4-cutoff}，我们注意到 $\mu_F$ 的二次项是虚假的，因为它唯一的用途就是抵消系数函数中的类似贡献。因此它对任何物理观测量都没有贡献。幂修正 renormalon 模型的思想 \cite{Beneke:1995pq,Dokshitzer:1995qm,Dasgupta:1996hh,Dasgupta:1996ki,Beneke:1997sr,Braun:2004bu} 是：把 $\mu_F$ 替换为一个适当的非微扰能标后，这一贡献反映了实际幂次压低贡献的阶数和函数形式。若接受这一``紫外主导''假设~\cite{Braun:1995tp,Beneke:1998ui,Beneke:2000kc}，就得到如下模型：
\begin{equation}
 \mathcal{Q}_4(x,p,\mu_F) = \kappa
\Lambda_{\rm QCD}^2 \int_{-1}^1 \!\frac{dy}{|y|} D_{\Qd}(\tfrac{x}{y})\, q(y,\mu_F)\,,
\label{UV_Lambda_int}
\end{equation}
其中无量纲系数 $\kappa = \mathcal{O}(1)$ 无法在理论内部确定，仍是一个自由参数。

在本工作中我们对不同版本的 qPDF（pPDF）在所谓的气泡链近似~\cite{Beneke:1998ui}下计算函数 $D_{\Qd}(x)$，这是我们的主要结果。该计算揭示，准分布与赝分布的幂修正具有如下普适结构
\begin{align}
 \cQ (x,p) &=  q(x) \Big\{1 + \mathcal{O}\Big( \frac{\Lambda^2}{p^2x^2(1-x)}\Big) \Big\},
\notag\\
\mathcal{P}(x,z) &= q(x)\Big\{1 + \mathcal{O}\big( z^2\!\Lambda^2{(1-x)}\big)\Big\},
\end{align}
并且会受到归一化的显著影响。特别地，把相关矩阵元按其在零动量处的取值归一化，会在较小 $x$ 处大幅减小 qPDF 的幂修正，但在较大 $x$ 处造成强烈放大。我们强调，当 $x\to 0$ 时 qPDF 的领头幂修正被 Bjorken $x$ 变量的两个幂次增强。对于 pPDF，幂修正在 $x\to1$ 处被压低；遗憾的是，经过文献~\cite{Orginos:2017kos} 的归一化步骤后这一性质不再成立（但换一种归一化即可保持）。此外，作为气泡链计算的副产品，我们还得到了微扰论各阶下领头扭度系数函数的大-$n_f$ 部分 $n^k_f \alpha_s^{k+1}$。



本文的结构安排如下。在第 2 节中我们更精确地表述我们的纲领：把 qPDF 和 pPDF 定义为位置空间（Ioffe 时间）分布的特定 Fourier 变换，简要讨论光线算符乘积展开（OPE）与靶质量修正，并介绍全文将要用到的相关技术（Borel 变换）和系统性近似（大-$n_f$ 展开）。第 3 节给出领头扭度系数函数 Borel 变换的结果并讨论其奇点结构。各类准分布的领头幂修放在第 4 节中得到。我们在那里汇集相关的解析表达式，并用价夸克 PDF 的现实参数化给出数值研究结果。最后的第 5 节用于总结与结论。


%
